\section{18. Doomsday}

Doomsday is the final escalation. The dungeon begins permanently destroying abandoned chambers. Priorities should shift from exploration and rewards toward survival and escape.

\subsection{Activation conditions}

Doomsday activates when a player enters treasure or the number of undiscovered rooms falls to four or fewer. Discovery is based on rooms revealed by opening doors, not physically entered. Merely opening the treasure door does not activate its entry trigger, though revealing it may independently satisfy the discovery condition.

\subsection{Activation timing}

Automatic activation occurs after the current door closes and after that closure's collapse checks. This distinction is essential.

1 — Enter treasure room

2 — Doomsday activation becomes pending

3 — Door closes

4 — Collapse checks use previous doomsday state

5 — DOOMSDAY ACTIVATES

6 — Future eligible closures cause collapses

The room from which the party enters treasure must not collapse solely because entering treasure activated doomsday. The same principle applies when the discovery threshold is reached. Manual activation is immediate but should normally be reserved for deliberate narrative interventions or corrections.

\subsection{Narrating activation}

\begin{quote}\itshape A deep vibration runs through the floor. The familiar grinding of moving stone is replaced by a thunderous cracking sound. Dust falls from the ceiling, and the walls tremble as though something enormous is tearing the dungeon apart.\end{quote}

\begin{quote}\itshape Whatever held this place together is failing.\end{quote}

Players need not understand every collapse restriction immediately, but should recognize that remaining has become significantly more dangerous.
